\documentclass[11pt]{article}

\usepackage[margin=0.85in]{geometry}
\usepackage{amsmath,amssymb}
\usepackage{booktabs}
\usepackage[hidelinks]{hyperref}

\newcommand{\argmin}{\operatorname*{arg\,min}}

\title{Three-Player, Three-Level Benchmark with\\ Inequalities as the Innermost Preference}
\author{}
\date{}

\begin{document}
\maketitle

This is the exact problem solved by
\texttt{test/quasi\_vs\_reduced\_preference\_inequality.jl}. It is built in two steps:

\begin{enumerate}
  \item The call
        \begin{center}
          \texttt{build\_benchmark\_problem(; num\_players = 3, levels = 3, kind)}
        \end{center}
        in \texttt{test/benchmark\_problems.jl} constructs a three-level generalized
        Nash problem with \emph{hard} coupled inequality constraints.
  \item The helper \texttt{innermost\_preference\_inequality\_goop} then deletes
        those hard constraints and re-attaches each player's inequality vector as a
        \emph{fourth, highest-priority preference level}, flagged with
        \begin{center}
          \texttt{is\_prioritized\_constraint = true}.
        \end{center}
        This is the encoding that \texttt{experiments/robotic\_arm\_core.jl} uses.
\end{enumerate}

Everything below is fully expanded --- no symbolic shorthand for the constants.

\section*{1. Notation}

There are three players. Player \(p \in \{1,2,3\}\) owns a decision vector
\[
  x_p = (x_{p,1},\, x_{p,2},\, x_{p,3},\, x_{p,4},\, x_{p,5}) \in \mathbb{R}^5,
  \qquad
  x = (x_1,x_2,x_3) \in \mathbb{R}^{15}.
\]
In code, \(x_{p,j}\) is \texttt{x[Block(p)][j]}. Write \(x_{-p}\) for the other two players'
vectors, which player \(p\) treats as fixed parameters --- the three player problems hold
simultaneously, i.e.\ this is a generalized Nash equilibrium problem. There are parameters
\(\theta \in \mathbb{R}^3\), but every objective and constraint below ignores them; they are
identically zero in the test.

\paragraph{Preference ordering.} Preferences are stored \texttt{[outermost, \dots, innermost]},
and the \emph{innermost is the highest priority}. Player \(p\) has four levels:

\begin{center}
\begin{tabular}{cllc}
  \toprule
  level \(\ell\) & controls & role & priority \\
  \midrule
  1 & \(x_{p,5}\)            & objective              & lowest \\
  2 & \(x_{p,3},\,x_{p,4}\)  & objective              & \\
  3 & \(x_{p,1},\,x_{p,2}\)  & objective              & \\
  4 & \(h_p(x) \ge 0\)       & prioritized constraint & \textbf{highest} \\
  \bottomrule
\end{tabular}
\end{center}

Each level's objective still ranges over the \emph{whole} vector \(x_p\); the coordinate lists
say which coordinates that level's objective actually depends on.

\section*{2. The prioritized-constraint level}

A level flagged \texttt{is\_prioritized\_constraint} does \emph{not} enter as a constraint. It
enters as the smooth penalty \texttt{ReducedGOOP.smooth\_piecewise\_preference\_objective},
applied elementwise and summed. In code that function is
\texttt{ifelse(h \(\ge\) 0, 0.0, (0 - h)\^{}(level + 2))}, i.e.
\[
  \varphi(h,\ell) =
  \begin{cases}
    0, & h \ge 0,\\
    (-h)^{\ell+2}, & h < 0.
  \end{cases}
\]
The inequality sits at \(\ell = 4\), so the exponent is \(\ell + 2 = 6\):
\[
  P_p(x) \;=\; \sum_{m=1}^{4} \varphi\bigl(h_{p,m}(x),\,4\bigr)
  \;=\; \sum_{m=1}^{4} \max\bigl(-h_{p,m}(x),\,0\bigr)^{6}.
\]
We have \(P_p \ge 0\) always, and \(P_p(x) = 0\) \emph{iff} every \(h_{p,m}(x) \ge 0\). So this
level's argmin set is exactly the feasible set of the original hard constraints, whenever that
set is nonempty.

\section*{3. The nested problem}

Player \(p\) solves the outermost level subject to being optimal for every higher-priority
level:
\[
\begin{aligned}
  \min_{x_p \in \mathbb{R}^5} \quad & J_p^{(1)}(x_p) \\
  \text{subject to} \quad & x_p \in \mathcal{S}_p^{(2)}(x_{-p}),
\end{aligned}
\]
where each level's solution set is defined by the next one down:
\[
\begin{aligned}
  \mathcal{S}_p^{(2)}(x_{-p}) &= \argmin_{x_p \in \mathbb{R}^5} \; J_p^{(2)}(x_p)
    \quad \text{subject to} \quad x_p \in \mathcal{S}_p^{(3)}(x_{-p}), \\
  \mathcal{S}_p^{(3)}(x_{-p}) &= \argmin_{x_p \in \mathbb{R}^5} \; J_p^{(3)}(x_p)
    \quad \text{subject to} \quad x_p \in \mathcal{S}_p^{(4)}(x_{-p}), \\
  \mathcal{S}_p^{(4)}(x_{-p}) &= \argmin_{x_p \in \mathbb{R}^5} \; P_p(x_p, x_{-p}).
\end{aligned}
\]
The innermost problem is \emph{unconstrained} --- that is the whole point of the
transformation. There are no hard constraints anywhere in this problem.

\section*{4. Coupling and constraint data}

The coupled resource consumed by player \(p\) in coordinate \(j\) is
\[
  c_{p,j}(x) \;=\; x_{p,j} \;+\; 0.25 \sum_{k \neq p} x_{k,j},
\]
with coupling weight \(0.25\) (\texttt{COUPLING\_WEIGHT}). Each player has four inequality rows:
\[
  h_p(x) =
  \begin{pmatrix}
    C_{p,1} - c_{p,1}(x) \\
    C_{p,2} - c_{p,2}(x) \\
    c_{p,3}(x) - F_{p,3} \\
    c_{p,4}(x) - F_{p,4}
  \end{pmatrix} \ge 0 .
\]
Rows 1 and 3 are \emph{active} at the designed solution; rows 2 and 4 carry slack \(1\).

\begin{center}
\begin{tabular}{ccccc}
  \toprule
  \(p\) & \(C_{p,1}\) & \(C_{p,2}\) & \(F_{p,3}\) & \(F_{p,4}\) \\
  \midrule
  1 & 0.6375 & 2.4250 & 0.8625 & 0.8750 \\
  2 & 0.6750 & 2.5000 & 0.9000 & 0.9500 \\
  3 & 0.7125 & 2.5750 & 0.9375 & 1.0250 \\
  \bottomrule
\end{tabular}
\qquad
\begin{tabular}{cccccc}
  \toprule
  \(p\) & \(t_{p,1}\) & \(t_{p,2}\) & \(t_{p,3}\) & \(t_{p,4}\) & \(t_{p,5}\) \\
  \midrule
  1 & 1.00 & 0.90 & 0.15 & 1.20 & 0.80 \\
  2 & 1.05 & 1.00 & 0.20 & 1.30 & 0.90 \\
  3 & 1.10 & 1.10 & 0.25 & 1.40 & 1.00 \\
  \bottomrule
\end{tabular}
\end{center}

\noindent
The right-hand table lists the unconstrained objective targets \(t_p\) --- the point each
objective would pick with no constraints.

\section*{5. \texttt{kind = :quadratic}, fully expanded}

Objectives are squared distances to the targets; the constraint rows are affine.

\subsection*{Player 1}
\[
\begin{aligned}
  \min_{x_1 \in \mathbb{R}^5} \quad & (x_{1,5}-0.80)^2 \\
  \text{subject to} \quad & x_1 \in \mathcal{S}_1^{(2)}(x_2,x_3)
\end{aligned}
\]
\[
\begin{aligned}
  \mathcal{S}_1^{(2)} &= \argmin_{x_1} \;
    (x_{1,3}-0.15)^2 + (x_{1,4}-1.20)^2
    \quad \text{s.t.} \quad x_1 \in \mathcal{S}_1^{(3)}, \\
  \mathcal{S}_1^{(3)} &= \argmin_{x_1} \;
    (x_{1,1}-1.00)^2 + (x_{1,2}-0.90)^2
    \quad \text{s.t.} \quad x_1 \in \mathcal{S}_1^{(4)}, \\
  \mathcal{S}_1^{(4)} &= \argmin_{x_1} \;
    \Bigl[ \max\bigl(-(0.6375 - x_{1,1} - 0.25x_{2,1} - 0.25x_{3,1}),\,0\bigr)^6 \\
    &\qquad\qquad + \max\bigl(-(2.4250 - x_{1,2} - 0.25x_{2,2} - 0.25x_{3,2}),\,0\bigr)^6 \\
    &\qquad\qquad + \max\bigl(-(x_{1,3} + 0.25x_{2,3} + 0.25x_{3,3} - 0.8625),\,0\bigr)^6 \\
    &\qquad\qquad + \max\bigl(-(x_{1,4} + 0.25x_{2,4} + 0.25x_{3,4} - 0.8750),\,0\bigr)^6 \Bigr].
\end{aligned}
\]

\subsection*{Player 2}
\[
\begin{aligned}
  \min_{x_2 \in \mathbb{R}^5} \quad & (x_{2,5}-0.90)^2 \\
  \text{subject to} \quad & x_2 \in \mathcal{S}_2^{(2)}(x_1,x_3)
\end{aligned}
\]
\[
\begin{aligned}
  \mathcal{S}_2^{(2)} &= \argmin_{x_2} \;
    (x_{2,3}-0.20)^2 + (x_{2,4}-1.30)^2
    \quad \text{s.t.} \quad x_2 \in \mathcal{S}_2^{(3)}, \\
  \mathcal{S}_2^{(3)} &= \argmin_{x_2} \;
    (x_{2,1}-1.05)^2 + (x_{2,2}-1.00)^2
    \quad \text{s.t.} \quad x_2 \in \mathcal{S}_2^{(4)}, \\
  \mathcal{S}_2^{(4)} &= \argmin_{x_2} \;
    \Bigl[ \max\bigl(-(0.6750 - x_{2,1} - 0.25x_{1,1} - 0.25x_{3,1}),\,0\bigr)^6 \\
    &\qquad\qquad + \max\bigl(-(2.5000 - x_{2,2} - 0.25x_{1,2} - 0.25x_{3,2}),\,0\bigr)^6 \\
    &\qquad\qquad + \max\bigl(-(x_{2,3} + 0.25x_{1,3} + 0.25x_{3,3} - 0.9000),\,0\bigr)^6 \\
    &\qquad\qquad + \max\bigl(-(x_{2,4} + 0.25x_{1,4} + 0.25x_{3,4} - 0.9500),\,0\bigr)^6 \Bigr].
\end{aligned}
\]

\subsection*{Player 3}
\[
\begin{aligned}
  \min_{x_3 \in \mathbb{R}^5} \quad & (x_{3,5}-1.00)^2 \\
  \text{subject to} \quad & x_3 \in \mathcal{S}_3^{(2)}(x_1,x_2)
\end{aligned}
\]
\[
\begin{aligned}
  \mathcal{S}_3^{(2)} &= \argmin_{x_3} \;
    (x_{3,3}-0.25)^2 + (x_{3,4}-1.40)^2
    \quad \text{s.t.} \quad x_3 \in \mathcal{S}_3^{(3)}, \\
  \mathcal{S}_3^{(3)} &= \argmin_{x_3} \;
    (x_{3,1}-1.10)^2 + (x_{3,2}-1.10)^2
    \quad \text{s.t.} \quad x_3 \in \mathcal{S}_3^{(4)}, \\
  \mathcal{S}_3^{(4)} &= \argmin_{x_3} \;
    \Bigl[ \max\bigl(-(0.7125 - x_{3,1} - 0.25x_{1,1} - 0.25x_{2,1}),\,0\bigr)^6 \\
    &\qquad\qquad + \max\bigl(-(2.5750 - x_{3,2} - 0.25x_{1,2} - 0.25x_{2,2}),\,0\bigr)^6 \\
    &\qquad\qquad + \max\bigl(-(x_{3,3} + 0.25x_{1,3} + 0.25x_{2,3} - 0.9375),\,0\bigr)^6 \\
    &\qquad\qquad + \max\bigl(-(x_{3,4} + 0.25x_{1,4} + 0.25x_{2,4} - 1.0250),\,0\bigr)^6 \Bigr].
\end{aligned}
\]

\section*{6. \texttt{kind = :nonlinear}, fully expanded}

Same structure, two substitutions:
\begin{itemize}
  \item every squared term \((u-t)^2\) becomes \(\cosh(u-t)-1\);
  \item every constraint row \(h\) becomes \(e^{h}-1\).
\end{itemize}
The second substitution preserves the feasible set and the active set exactly, since
\(e^{r}-1 = 0 \iff r = 0\) and \(e^{r}-1 > 0 \iff r > 0\). It does change the \emph{geometry}:
the rows are no longer affine, and \(\cosh\) has nonvanishing derivatives of every order. Note
\(-(e^{h}-1) = 1-e^{h}\), which is how the penalty is written below.

\subsection*{Player 1}
\[
\begin{aligned}
  \min_{x_1 \in \mathbb{R}^5} \quad & \cosh(x_{1,5}-0.80)-1 \\
  \text{subject to} \quad & x_1 \in \mathcal{S}_1^{(2)}(x_2,x_3)
\end{aligned}
\]
\[
\begin{aligned}
  \mathcal{S}_1^{(2)} &= \argmin_{x_1} \;
    \bigl[\cosh(x_{1,3}-0.15)-1\bigr] + \bigl[\cosh(x_{1,4}-1.20)-1\bigr]
    \quad \text{s.t.} \quad x_1 \in \mathcal{S}_1^{(3)}, \\
  \mathcal{S}_1^{(3)} &= \argmin_{x_1} \;
    \bigl[\cosh(x_{1,1}-1.00)-1\bigr] + \bigl[\cosh(x_{1,2}-0.90)-1\bigr]
    \quad \text{s.t.} \quad x_1 \in \mathcal{S}_1^{(4)}, \\
  \mathcal{S}_1^{(4)} &= \argmin_{x_1} \;
    \Bigl[ \max\bigl(1 - e^{\,0.6375 - x_{1,1} - 0.25x_{2,1} - 0.25x_{3,1}},\,0\bigr)^6 \\
    &\qquad\qquad + \max\bigl(1 - e^{\,2.4250 - x_{1,2} - 0.25x_{2,2} - 0.25x_{3,2}},\,0\bigr)^6 \\
    &\qquad\qquad + \max\bigl(1 - e^{\,x_{1,3} + 0.25x_{2,3} + 0.25x_{3,3} - 0.8625},\,0\bigr)^6 \\
    &\qquad\qquad + \max\bigl(1 - e^{\,x_{1,4} + 0.25x_{2,4} + 0.25x_{3,4} - 0.8750},\,0\bigr)^6 \Bigr].
\end{aligned}
\]

\subsection*{Player 2}
\[
\begin{aligned}
  \min_{x_2 \in \mathbb{R}^5} \quad & \cosh(x_{2,5}-0.90)-1 \\
  \text{subject to} \quad & x_2 \in \mathcal{S}_2^{(2)}(x_1,x_3)
\end{aligned}
\]
\[
\begin{aligned}
  \mathcal{S}_2^{(2)} &= \argmin_{x_2} \;
    \bigl[\cosh(x_{2,3}-0.20)-1\bigr] + \bigl[\cosh(x_{2,4}-1.30)-1\bigr]
    \quad \text{s.t.} \quad x_2 \in \mathcal{S}_2^{(3)}, \\
  \mathcal{S}_2^{(3)} &= \argmin_{x_2} \;
    \bigl[\cosh(x_{2,1}-1.05)-1\bigr] + \bigl[\cosh(x_{2,2}-1.00)-1\bigr]
    \quad \text{s.t.} \quad x_2 \in \mathcal{S}_2^{(4)}, \\
  \mathcal{S}_2^{(4)} &= \argmin_{x_2} \;
    \Bigl[ \max\bigl(1 - e^{\,0.6750 - x_{2,1} - 0.25x_{1,1} - 0.25x_{3,1}},\,0\bigr)^6 \\
    &\qquad\qquad + \max\bigl(1 - e^{\,2.5000 - x_{2,2} - 0.25x_{1,2} - 0.25x_{3,2}},\,0\bigr)^6 \\
    &\qquad\qquad + \max\bigl(1 - e^{\,x_{2,3} + 0.25x_{1,3} + 0.25x_{3,3} - 0.9000},\,0\bigr)^6 \\
    &\qquad\qquad + \max\bigl(1 - e^{\,x_{2,4} + 0.25x_{1,4} + 0.25x_{3,4} - 0.9500},\,0\bigr)^6 \Bigr].
\end{aligned}
\]

\subsection*{Player 3}
\[
\begin{aligned}
  \min_{x_3 \in \mathbb{R}^5} \quad & \cosh(x_{3,5}-1.00)-1 \\
  \text{subject to} \quad & x_3 \in \mathcal{S}_3^{(2)}(x_1,x_2)
\end{aligned}
\]
\[
\begin{aligned}
  \mathcal{S}_3^{(2)} &= \argmin_{x_3} \;
    \bigl[\cosh(x_{3,3}-0.25)-1\bigr] + \bigl[\cosh(x_{3,4}-1.40)-1\bigr]
    \quad \text{s.t.} \quad x_3 \in \mathcal{S}_3^{(3)}, \\
  \mathcal{S}_3^{(3)} &= \argmin_{x_3} \;
    \bigl[\cosh(x_{3,1}-1.10)-1\bigr] + \bigl[\cosh(x_{3,2}-1.10)-1\bigr]
    \quad \text{s.t.} \quad x_3 \in \mathcal{S}_3^{(4)}, \\
  \mathcal{S}_3^{(4)} &= \argmin_{x_3} \;
    \Bigl[ \max\bigl(1 - e^{\,0.7125 - x_{3,1} - 0.25x_{1,1} - 0.25x_{2,1}},\,0\bigr)^6 \\
    &\qquad\qquad + \max\bigl(1 - e^{\,2.5750 - x_{3,2} - 0.25x_{1,2} - 0.25x_{2,2}},\,0\bigr)^6 \\
    &\qquad\qquad + \max\bigl(1 - e^{\,x_{3,3} + 0.25x_{1,3} + 0.25x_{2,3} - 0.9375},\,0\bigr)^6 \\
    &\qquad\qquad + \max\bigl(1 - e^{\,x_{3,4} + 0.25x_{1,4} + 0.25x_{2,4} - 1.0250},\,0\bigr)^6 \Bigr].
\end{aligned}
\]

\section*{7. Ground truth}

The family is constructed so that the hard-constrained problem has the closed-form solution
\[
  x_1^\star = \begin{pmatrix}0.40\\0.90\\0.55\\1.20\\0.80\end{pmatrix},\qquad
  x_2^\star = \begin{pmatrix}0.45\\1.00\\0.60\\1.30\\0.90\end{pmatrix},\qquad
  x_3^\star = \begin{pmatrix}0.50\\1.10\\0.65\\1.40\\1.00\end{pmatrix}.
\]
\texttt{case.expected} in the test is this vector, flattened to length 15.

\paragraph{How it is constructed.} \(t_p\) differs from \(x_p^\star\) only in coordinates 1
and 3:
\[
  t_{p,1} = x_{p,1}^\star + 0.6, \qquad t_{p,3} = x_{p,3}^\star - 0.4,
\]
with \(t_{p,2} = x_{p,2}^\star\), \(t_{p,4} = x_{p,4}^\star\), \(t_{p,5} = x_{p,5}^\star\).
Coordinates 2, 4 and 5 are therefore already at their unconstrained optimum and no constraint
binds them. Coordinates 1 and 3 are pulled off target by exactly the two active rows:
\(C_{p,1}\) is \emph{defined} as \(c_{p,1}(x^\star)\) and \(F_{p,3}\) as \(c_{p,3}(x^\star)\),
while the two inactive rows are offset by \(\pm 1\).

\paragraph{Verification at \(x^\star\)} (identical for all three players):

\begin{center}
\begin{tabular}{cccl}
  \toprule
  row & \(h_{p,m}(x^\star)\), \texttt{:quadratic} & \(e^{h}-1\), \texttt{:nonlinear} & status \\
  \midrule
  1 & \(0\) & \(0\)                    & active \\
  2 & \(1\) & \(e-1 \approx 1.71828\)  & inactive \\
  3 & \(0\) & \(0\)                    & active \\
  4 & \(1\) & \(e-1 \approx 1.71828\)  & inactive \\
  \bottomrule
\end{tabular}
\end{center}

Hence \(P_p(x^\star) = 0\) for every player: \(x^\star\) attains the exact minimum of the
highest-priority level. Since that level's argmin set equals the feasible set, the nested
problem above has the \emph{same} solution as the hard-constrained original --- in exact
arithmetic.

The \(3\times 3\) coupling matrix has \(1\) on the diagonal and \(0.25\) off it, so it is
nonsingular and the two active rows pin each player's share individually, not just the
aggregate.

\section*{8. Why solves do not land on \texorpdfstring{\(x^\star\)}{x*}}

\(x^\star\) sits exactly on the boundary of the two active rows, where
\(h_{p,1} = h_{p,3} = 0\). The penalty \(\max(-h,0)^6\) has value \emph{and first five
derivatives all zero} there. The highest-priority level therefore exerts no gradient at the
solution: nothing pulls an iterate back onto the active constraints, and the lower-priority
objectives --- which want coordinate 1 larger and coordinate 3 smaller --- push it into mild
infeasibility until the sixth power grows enough to resist.

Measured with the test's configuration (KLU, backtracking, \(\text{tol} = 10^{-4}\),
\texttt{max\_outer\_iters = 1}):

\begin{center}
\begin{tabular}{cllc}
  \toprule
  \texttt{kind} & \(z_0\) & \(\lVert x - x^\star\rVert_\infty\) reduced / quasi & max violation \\
  \midrule
  \texttt{:quadratic} & \texttt{expected} & 0.06068 / 0.06068 & 0.0910 \\
  \texttt{:quadratic} & \texttt{zeros}    & 0.72728 / 0.72728 & 0.0909 \\
  \texttt{:quadratic} & \texttt{offset}   & 0.06057 / 0.06057 & 0.0909 \\
  \texttt{:nonlinear} & \texttt{expected} & 0.06257 / 0.06257 & 0.0896 \\
  \texttt{:nonlinear} & \texttt{zeros}    & 0.73006 / 0.73006 & 0.0907 \\
  \texttt{:nonlinear} & \texttt{offset}   & 0.06240 / 0.06240 & 0.0894 \\
  \bottomrule
\end{tabular}
\end{center}

The residual is not what limits accuracy here --- the degenerate penalty is. Tightening the
tolerance to \(10^{-6}\) makes every solve report \texttt{:failed} (the iteration stalls at
\(\lVert F\rVert \approx 1.5\times 10^{-5} \dots 2.9\times 10^{-5}\)) yet moves the answer
only from \(0.061\) to \(0.046\) away from \(x^\star\), with the violation shrinking from
\(0.091\) to \(0.069\).

This gap is a property of the \((-h)^6\) encoding, not of the \texttt{:reduced} /
\texttt{:quasi} choice: the two formulations agree with each other to between
\(3\times 10^{-15}\) and \(4\times 10^{-6}\), i.e.\ four to nine orders of magnitude closer to
each other than either is to \(x^\star\).

It is the same encoding \texttt{robotic\_arm\_core.jl} uses for its collision and speed
constraints, so the same slack should be expected there --- a violation of \(10^{-2}\)
registers as a penalty of only \(10^{-12}\).

\end{document}
